% Options for packages loaded elsewhere
\PassOptionsToPackage{unicode}{hyperref}
\PassOptionsToPackage{hyphens}{url}
\PassOptionsToPackage{dvipsnames,svgnames,x11names}{xcolor}
\documentclass[
  11pt,
  letterpaper,
  titlepage]{article}
\usepackage{xcolor}
\usepackage[margin=0.85in]{geometry}
\usepackage{amsmath,amssymb}
\setcounter{secnumdepth}{5}
\usepackage{iftex}
\ifPDFTeX
  \usepackage[T1]{fontenc}
  \usepackage[utf8]{inputenc}
  \usepackage{textcomp} % provide euro and other symbols
\else % if luatex or xetex
  \usepackage{unicode-math} % this also loads fontspec
  \defaultfontfeatures{Scale=MatchLowercase}
  \defaultfontfeatures[\rmfamily]{Ligatures=TeX,Scale=1}
\fi
\usepackage{lmodern}
\ifPDFTeX\else
  % xetex/luatex font selection
\fi
% Use upquote if available, for straight quotes in verbatim environments
\IfFileExists{upquote.sty}{\usepackage{upquote}}{}
\IfFileExists{microtype.sty}{% use microtype if available
  \usepackage[]{microtype}
  \UseMicrotypeSet[protrusion]{basicmath} % disable protrusion for tt fonts
}{}
\usepackage{setspace}
\makeatletter
\@ifundefined{KOMAClassName}{% if non-KOMA class
  \IfFileExists{parskip.sty}{%
    \usepackage{parskip}
  }{% else
    \setlength{\parindent}{0pt}
    \setlength{\parskip}{6pt plus 2pt minus 1pt}}
}{% if KOMA class
  \KOMAoptions{parskip=half}}
\makeatother
\usepackage{color}
\usepackage{fancyvrb}
\newcommand{\VerbBar}{|}
\newcommand{\VERB}{\Verb[commandchars=\\\{\}]}
\DefineVerbatimEnvironment{Highlighting}{Verbatim}{commandchars=\\\{\}}
% Add ',fontsize=\small' for more characters per line
\usepackage{framed}
\definecolor{shadecolor}{RGB}{248,248,248}
\newenvironment{Shaded}{\begin{snugshade}}{\end{snugshade}}
\newcommand{\AlertTok}[1]{\textcolor[rgb]{0.94,0.16,0.16}{#1}}
\newcommand{\AnnotationTok}[1]{\textcolor[rgb]{0.56,0.35,0.01}{\textbf{\textit{#1}}}}
\newcommand{\AttributeTok}[1]{\textcolor[rgb]{0.13,0.29,0.53}{#1}}
\newcommand{\BaseNTok}[1]{\textcolor[rgb]{0.00,0.00,0.81}{#1}}
\newcommand{\BuiltInTok}[1]{#1}
\newcommand{\CharTok}[1]{\textcolor[rgb]{0.31,0.60,0.02}{#1}}
\newcommand{\CommentTok}[1]{\textcolor[rgb]{0.56,0.35,0.01}{\textit{#1}}}
\newcommand{\CommentVarTok}[1]{\textcolor[rgb]{0.56,0.35,0.01}{\textbf{\textit{#1}}}}
\newcommand{\ConstantTok}[1]{\textcolor[rgb]{0.56,0.35,0.01}{#1}}
\newcommand{\ControlFlowTok}[1]{\textcolor[rgb]{0.13,0.29,0.53}{\textbf{#1}}}
\newcommand{\DataTypeTok}[1]{\textcolor[rgb]{0.13,0.29,0.53}{#1}}
\newcommand{\DecValTok}[1]{\textcolor[rgb]{0.00,0.00,0.81}{#1}}
\newcommand{\DocumentationTok}[1]{\textcolor[rgb]{0.56,0.35,0.01}{\textbf{\textit{#1}}}}
\newcommand{\ErrorTok}[1]{\textcolor[rgb]{0.64,0.00,0.00}{\textbf{#1}}}
\newcommand{\ExtensionTok}[1]{#1}
\newcommand{\FloatTok}[1]{\textcolor[rgb]{0.00,0.00,0.81}{#1}}
\newcommand{\FunctionTok}[1]{\textcolor[rgb]{0.13,0.29,0.53}{\textbf{#1}}}
\newcommand{\ImportTok}[1]{#1}
\newcommand{\InformationTok}[1]{\textcolor[rgb]{0.56,0.35,0.01}{\textbf{\textit{#1}}}}
\newcommand{\KeywordTok}[1]{\textcolor[rgb]{0.13,0.29,0.53}{\textbf{#1}}}
\newcommand{\NormalTok}[1]{#1}
\newcommand{\OperatorTok}[1]{\textcolor[rgb]{0.81,0.36,0.00}{\textbf{#1}}}
\newcommand{\OtherTok}[1]{\textcolor[rgb]{0.56,0.35,0.01}{#1}}
\newcommand{\PreprocessorTok}[1]{\textcolor[rgb]{0.56,0.35,0.01}{\textit{#1}}}
\newcommand{\RegionMarkerTok}[1]{#1}
\newcommand{\SpecialCharTok}[1]{\textcolor[rgb]{0.81,0.36,0.00}{\textbf{#1}}}
\newcommand{\SpecialStringTok}[1]{\textcolor[rgb]{0.31,0.60,0.02}{#1}}
\newcommand{\StringTok}[1]{\textcolor[rgb]{0.31,0.60,0.02}{#1}}
\newcommand{\VariableTok}[1]{\textcolor[rgb]{0.00,0.00,0.00}{#1}}
\newcommand{\VerbatimStringTok}[1]{\textcolor[rgb]{0.31,0.60,0.02}{#1}}
\newcommand{\WarningTok}[1]{\textcolor[rgb]{0.56,0.35,0.01}{\textbf{\textit{#1}}}}
\usepackage{longtable,booktabs,array}
\usepackage{caption}
\captionsetup[table]{skip=6pt}
\newcounter{none} % for unnumbered tables
\usepackage{calc} % for calculating minipage widths
% Correct order of tables after \paragraph or \subparagraph
\usepackage{etoolbox}
\makeatletter
\patchcmd\longtable{\par}{\if@noskipsec\mbox{}\fi\par}{}{}
\makeatother
% Allow footnotes in longtable head/foot
\IfFileExists{footnotehyper.sty}{\usepackage{footnotehyper}}{\usepackage{footnote}}
\makesavenoteenv{longtable}
\setlength{\emergencystretch}{3em} % prevent overfull lines
\providecommand{\tightlist}{%
  \setlength{\itemsep}{0pt}\setlength{\parskip}{0pt}}
\usepackage{bookmark}
\IfFileExists{xurl.sty}{\usepackage{xurl}}{} % add URL line breaks if available
\urlstyle{same}
% fallback for those not using the hyperref driver hyperxmp:
\makeatletter
\@ifundefined{xmpquote}{\newcommand{\xmpquote}[1]{#1}}{}
\makeatother
\hypersetup{
  pdftitle={Config Lib Integration Guide},
  colorlinks=true,
  linkcolor={blue},
  filecolor={Maroon},
  citecolor={Blue},
  urlcolor={blue},
  pdfcreator={LaTeX via pandoc}}

\title{Config Lib Integration Guide}
\usepackage{etoolbox}
\makeatletter
\providecommand{\subtitle}[1]{% add subtitle to \maketitle
  \apptocmd{\@title}{\par {\large #1 \par}}{}{}
}
\makeatother
\subtitle{Config Lib 1.12.0 for Vintage Story 1.22.6}
\author{}
\date{2026-08-15}

\begin{document}
\maketitle

{
\hypersetup{linkcolor=black}
\setcounter{tocdepth}{2}
\tableofcontents
}
\setstretch{1.08}
This document explains how a Vintage Story mod can support Config Lib
while continuing to work when Config Lib is not installed. It covers
content-mod asset patches, ordinary code-mod JSON configuration, live
setting updates, and the public runtime API visible in Config Lib
\texttt{1.12.0}.

The findings apply to the Config Lib \texttt{1.12.0} reference retained
for Vintage Story \texttt{1.22.6}. That Config Lib release was published
for the Vintage Story \texttt{1.22.x} line through \texttt{1.22.1}, so
the repository classifies it as a same-line fallback rather than an
exact \texttt{1.22.6} release. Recheck Config Lib when targeting another
game or Config Lib version.

\section{Short Answer}\label{short-answer}

Yes, a mod can implement Config Lib without making it a mandatory
runtime dependency.

The recommended design for an ordinary code mod is:

\begin{enumerate}
\def\labelenumi{\arabic{enumi}.}
\tightlist
\item
  Keep a normal configuration class and load it with Vintage Story's
  \texttt{LoadModConfig\textless{}T\textgreater{}()} API.
\item
  Supply sensible defaults that work with no Config Lib installation.
\item
  Add
  \texttt{assets/\textless{}modid\textgreater{}/config/configlib-patches.json}
  to describe the Config Lib controls.
\item
  Set \texttt{file} in that asset to the same JSON filename loaded by
  the mod.
\item
  If settings must take effect immediately, test
  \texttt{IsModEnabled("configlib")} before resolving
  \texttt{ConfigLibModSystem}, then subscribe to its configuration
  events.
\item
  Keep the defaults in the C\# class, assets, and Config Lib definition
  identical.
\end{enumerate}

Config Lib's upstream JSON-API documentation explicitly treats it as
optional. If it is absent, its configuration files and asset patches are
not applied; the consuming mod's ordinary JSON file or asset defaults
remain in effect.

\section{Choose An Integration Model}\label{choose-an-integration-model}

{\def\LTcaptype{none} % do not increment counter
\begin{longtable}[]{@{}
  >{\raggedright\arraybackslash}p{(\linewidth - 6\tabcolsep) * \real{0.2500}}
  >{\raggedright\arraybackslash}p{(\linewidth - 6\tabcolsep) * \real{0.2500}}
  >{\raggedright\arraybackslash}p{(\linewidth - 6\tabcolsep) * \real{0.2500}}
  >{\raggedright\arraybackslash}p{(\linewidth - 6\tabcolsep) * \real{0.2500}}@{}}
\toprule\noalign{}
\begin{minipage}[b]{\linewidth}\raggedright
Model
\end{minipage} & \begin{minipage}[b]{\linewidth}\raggedright
Saved File
\end{minipage} & \begin{minipage}[b]{\linewidth}\raggedright
Best For
\end{minipage} & \begin{minipage}[b]{\linewidth}\raggedright
C\# API Needed
\end{minipage} \\
\midrule\noalign{}
\endhead
\bottomrule\noalign{}
\endlastfoot
Generated YAML & Domain YAML file & Content mods and asset patching &
No \\
JSON-backed asset & Named JSON file & Existing code-mod config classes &
Only for live callbacks \\
Managed registration & JSON path from code & Code-first configuration &
Yes \\
\end{longtable}
}

The first two models are discovered from
\texttt{configlib-patches.json}. The third is registered through Config
Lib's public runtime API and must be registered early enough that Config
Lib has not already sent its configuration registry to clients.

\section{The Required Asset Location}\label{the-required-asset-location}

Config Lib \texttt{1.12.0} scans the Vintage Story \texttt{config} asset
category and selects assets whose filename is exactly:

\begin{Shaded}
\begin{Highlighting}[]
\NormalTok{configlib{-}patches.json}
\end{Highlighting}
\end{Shaded}

Place the file at:

\begin{Shaded}
\begin{Highlighting}[]
\NormalTok{assets/\textless{}modid\textgreater{}/config/configlib{-}patches.json}
\end{Highlighting}
\end{Shaded}

The asset domain becomes the Config Lib configuration domain. For a mod
id of \texttt{examplemod}, the configuration is therefore retrieved with
\texttt{GetConfig("examplemod")}, and Config Lib events report
\texttt{examplemod} as their domain.

The filename is significant. A differently named JSON file in the same
directory is not discovered by the loader.

\section{Recommended Code-Mod
Implementation}\label{recommended-code-mod-implementation}

\subsection{Keep A Normal Fallback
Config}\label{keep-a-normal-fallback-config}

Config Lib should enhance the mod's configuration experience rather than
own the only viable configuration path. A minimal ordinary config class
can look like this:

\begin{Shaded}
\begin{Highlighting}[]
\KeywordTok{public} \KeywordTok{sealed} \KeywordTok{class}\NormalTok{ ModConfig}
\OperatorTok{\{}
    \KeywordTok{public} \DataTypeTok{bool}\NormalTok{ Enabled }\OperatorTok{\{} \KeywordTok{get}\OperatorTok{;} \KeywordTok{set}\OperatorTok{;} \OperatorTok{\}} \OperatorTok{=} \KeywordTok{true}\OperatorTok{;}

    \KeywordTok{public} \DataTypeTok{int}\NormalTok{ SearchRadius }\OperatorTok{\{} \KeywordTok{get}\OperatorTok{;} \KeywordTok{set}\OperatorTok{;} \OperatorTok{\}} \OperatorTok{=} \DecValTok{5}\OperatorTok{;}

    \KeywordTok{public} \DataTypeTok{string}\NormalTok{ DisplayMode }\OperatorTok{\{} \KeywordTok{get}\OperatorTok{;} \KeywordTok{set}\OperatorTok{;} \OperatorTok{\}} \OperatorTok{=} \StringTok{"Detailed"}\OperatorTok{;}
\OperatorTok{\}}
\end{Highlighting}
\end{Shaded}

Load and store it through the normal Vintage Story API before attempting
Config Lib integration:

\begin{Shaded}
\begin{Highlighting}[]
\KeywordTok{private} \DataTypeTok{const} \DataTypeTok{string}\NormalTok{ ConfigFileName }\OperatorTok{=} \StringTok{"ExampleMod.json"}\OperatorTok{;}

\KeywordTok{public}\NormalTok{ ModConfig Config }\OperatorTok{\{} \KeywordTok{get}\OperatorTok{;} \KeywordTok{private} \KeywordTok{set}\OperatorTok{;} \OperatorTok{\}} \OperatorTok{=} \KeywordTok{new}\OperatorTok{();}

\KeywordTok{public} \KeywordTok{override} \DataTypeTok{void} \FunctionTok{Start}\OperatorTok{(}\NormalTok{ICoreAPI api}\OperatorTok{)}
\OperatorTok{\{}
    \KeywordTok{base}\OperatorTok{.}\FunctionTok{Start}\OperatorTok{(}\NormalTok{api}\OperatorTok{);}

    \KeywordTok{try}
    \OperatorTok{\{}
\NormalTok{        Config }\OperatorTok{=}\NormalTok{ api}\OperatorTok{.}\FunctionTok{LoadModConfig}\OperatorTok{\textless{}}\NormalTok{ModConfig}\OperatorTok{\textgreater{}(}\NormalTok{ConfigFileName}\OperatorTok{)} \OperatorTok{??} \KeywordTok{new} \FunctionTok{ModConfig}\OperatorTok{();}
\NormalTok{        api}\OperatorTok{.}\FunctionTok{StoreModConfig}\OperatorTok{(}\NormalTok{Config}\OperatorTok{,}\NormalTok{ ConfigFileName}\OperatorTok{);}
    \OperatorTok{\}}
    \KeywordTok{catch} \OperatorTok{(}\NormalTok{Exception exception}\OperatorTok{)}
    \OperatorTok{\{}
\NormalTok{        api}\OperatorTok{.}\FunctionTok{Logger}\OperatorTok{.}\FunctionTok{Error}\OperatorTok{(}\StringTok{"Unable to load ExampleMod configuration: \{0\}"}\OperatorTok{,}\NormalTok{ exception}\OperatorTok{);}
\NormalTok{        Config }\OperatorTok{=} \KeywordTok{new} \FunctionTok{ModConfig}\OperatorTok{();}
    \OperatorTok{\}}

    \FunctionTok{AttachOptionalConfigLib}\OperatorTok{(}\NormalTok{api}\OperatorTok{);}
\OperatorTok{\}}
\end{Highlighting}
\end{Shaded}

This fallback is the behavior that protects the mod when Config Lib is
missing, disabled, or unable to parse its definition. The exact
exception policy is project-specific, but the mod should always retain a
usable configuration object.

\subsection{Describe The Same JSON File To Config
Lib}\label{describe-the-same-json-file-to-config-lib}

The corresponding \texttt{configlib-patches.json} can use Config Lib
\texttt{1.12.0}'s array form:

\begin{Shaded}
\begin{Highlighting}[]
\FunctionTok{\{}
  \DataTypeTok{"version"}\FunctionTok{:} \DecValTok{1}\FunctionTok{,}
  \DataTypeTok{"file"}\FunctionTok{:} \StringTok{"ExampleMod.json"}\FunctionTok{,}
  \DataTypeTok{"settings"}\FunctionTok{:} \OtherTok{[}
    \FunctionTok{\{}
      \DataTypeTok{"code"}\FunctionTok{:} \StringTok{"Enabled"}\FunctionTok{,}
      \DataTypeTok{"type"}\FunctionTok{:} \StringTok{"boolean"}\FunctionTok{,}
      \DataTypeTok{"ingui"}\FunctionTok{:} \StringTok{"examplemod:setting{-}enabled"}\FunctionTok{,}
      \DataTypeTok{"comment"}\FunctionTok{:} \StringTok{"examplemod:setting{-}enabled{-}description"}\FunctionTok{,}
      \DataTypeTok{"default"}\FunctionTok{:} \KeywordTok{true}\FunctionTok{,}
      \DataTypeTok{"weight"}\FunctionTok{:} \DecValTok{1}
    \FunctionTok{\}}\OtherTok{,}
    \FunctionTok{\{}
      \DataTypeTok{"code"}\FunctionTok{:} \StringTok{"SearchRadius"}\FunctionTok{,}
      \DataTypeTok{"type"}\FunctionTok{:} \StringTok{"integer"}\FunctionTok{,}
      \DataTypeTok{"ingui"}\FunctionTok{:} \StringTok{"examplemod:setting{-}search{-}radius"}\FunctionTok{,}
      \DataTypeTok{"comment"}\FunctionTok{:} \StringTok{"examplemod:setting{-}search{-}radius{-}description"}\FunctionTok{,}
      \DataTypeTok{"default"}\FunctionTok{:} \DecValTok{5}\FunctionTok{,}
      \DataTypeTok{"range"}\FunctionTok{:} \FunctionTok{\{}
        \DataTypeTok{"min"}\FunctionTok{:} \DecValTok{1}\FunctionTok{,}
        \DataTypeTok{"max"}\FunctionTok{:} \DecValTok{20}\FunctionTok{,}
        \DataTypeTok{"step"}\FunctionTok{:} \DecValTok{1}
      \FunctionTok{\},}
      \DataTypeTok{"weight"}\FunctionTok{:} \DecValTok{2}
    \FunctionTok{\}}\OtherTok{,}
    \FunctionTok{\{}
      \DataTypeTok{"code"}\FunctionTok{:} \StringTok{"DisplayMode"}\FunctionTok{,}
      \DataTypeTok{"type"}\FunctionTok{:} \StringTok{"string"}\FunctionTok{,}
      \DataTypeTok{"ingui"}\FunctionTok{:} \StringTok{"examplemod:setting{-}display{-}mode"}\FunctionTok{,}
      \DataTypeTok{"default"}\FunctionTok{:} \StringTok{"Detailed"}\FunctionTok{,}
      \DataTypeTok{"values"}\FunctionTok{:} \OtherTok{[}
        \StringTok{"Compact"}\OtherTok{,}
        \StringTok{"Detailed"}
      \OtherTok{]}\FunctionTok{,}
      \DataTypeTok{"weight"}\FunctionTok{:} \DecValTok{3}
    \FunctionTok{\}}
  \OtherTok{]}
\FunctionTok{\}}
\end{Highlighting}
\end{Shaded}

When \texttt{file} is present, Config Lib treats the target as JSON
instead of generating its default
\texttt{ModConfig/\textless{}domain\textgreater{}.yaml} file. The
\texttt{code} values act as paths into that JSON document. Simple
top-level property names therefore match ordinary C\# property names
naturally.

Config Lib's JSON reader also uses its path implementation for nested
values. Test nested paths with the target version rather than assuming
that every JSONPath feature is available.

\subsection{Apply Initial And Live
Values}\label{apply-initial-and-live-values}

Code that names Config Lib types needs a compile-time reference to
Config Lib, but the runtime path can still be optional. The retained
Betterer Prospecting example first checks
\texttt{IsModEnabled("configlib")} and only then resolves
\texttt{ConfigLibModSystem} and subscribes to events.

A reusable pattern is:

\begin{Shaded}
\begin{Highlighting}[]
\KeywordTok{private} \DataTypeTok{void} \FunctionTok{AttachOptionalConfigLib}\OperatorTok{(}\NormalTok{ICoreAPI api}\OperatorTok{)}
\OperatorTok{\{}
    \KeywordTok{if} \OperatorTok{(!}\NormalTok{api}\OperatorTok{.}\FunctionTok{ModLoader}\OperatorTok{.}\FunctionTok{IsModEnabled}\OperatorTok{(}\StringTok{"configlib"}\OperatorTok{))}
    \OperatorTok{\{}
        \KeywordTok{return}\OperatorTok{;}
    \OperatorTok{\}}

\NormalTok{    ConfigLibModSystem provider }\OperatorTok{=}
\NormalTok{        api}\OperatorTok{.}\FunctionTok{ModLoader}\OperatorTok{.}\FunctionTok{GetModSystem}\OperatorTok{\textless{}}\NormalTok{ConfigLibModSystem}\OperatorTok{\textgreater{}(}\KeywordTok{true}\OperatorTok{);}

\NormalTok{    provider}\OperatorTok{.}\FunctionTok{ConfigsLoaded} \OperatorTok{+=} \OperatorTok{()} \OperatorTok{=\textgreater{}}
    \OperatorTok{\{}
\NormalTok{        IConfig}\OperatorTok{?}\NormalTok{ config }\OperatorTok{=}\NormalTok{ provider}\OperatorTok{.}\FunctionTok{GetConfig}\OperatorTok{(}\NormalTok{Mod}\OperatorTok{.}\FunctionTok{Info}\OperatorTok{.}\FunctionTok{ModID}\OperatorTok{);}
\NormalTok{        config}\OperatorTok{?.}\FunctionTok{AssignSettingsValues}\OperatorTok{(}\NormalTok{Config}\OperatorTok{);}
        \FunctionTok{ApplyConfigDependentState}\OperatorTok{();}
    \OperatorTok{\};}

\NormalTok{    provider}\OperatorTok{.}\FunctionTok{SettingChanged} \OperatorTok{+=} \OperatorTok{(}\NormalTok{domain}\OperatorTok{,}\NormalTok{ \_}\OperatorTok{,}\NormalTok{ setting}\OperatorTok{)} \OperatorTok{=\textgreater{}}
    \OperatorTok{\{}
        \KeywordTok{if} \OperatorTok{(}\NormalTok{domain }\OperatorTok{!=}\NormalTok{ Mod}\OperatorTok{.}\FunctionTok{Info}\OperatorTok{.}\FunctionTok{ModID}\OperatorTok{)}
        \OperatorTok{\{}
            \KeywordTok{return}\OperatorTok{;}
        \OperatorTok{\}}

\NormalTok{        setting}\OperatorTok{.}\FunctionTok{AssignSettingValue}\OperatorTok{(}\NormalTok{Config}\OperatorTok{);}
        \FunctionTok{ApplySettingChange}\OperatorTok{(}\NormalTok{setting}\OperatorTok{.}\FunctionTok{YamlCode}\OperatorTok{);}
    \OperatorTok{\};}
\OperatorTok{\}}
\end{Highlighting}
\end{Shaded}

\texttt{IConfig.AssignSettingsValues(object)} assigns all current
settings to matching fields or properties on an object.
\texttt{ISetting.AssignSettingValue(object)} assigns one setting. These
methods return or operate through Config Lib's current parsed values,
including values received from the server.

The callbacks that follow assignment are important when a setting
affects cached data, Harmony patches, registered tool modes, meshes, or
other state that was constructed earlier. Changing the config object's
property alone does not rebuild such state automatically.

Use the nullable syntax and event-unsubscription pattern appropriate to
the consuming project's target framework. Store delegate instances when
they must be removed during \texttt{Dispose()}; an inline lambda cannot
be reliably removed later by recreating a visually identical lambda.

\section{Setting Definition
Reference}\label{setting-definition-reference}

Config Lib \texttt{1.12.0} accepts an array of setting and formatting
blocks or the older grouped-object form documented by the upstream JSON
API.

\subsection{Array Form}\label{array-form}

Each setting entry normally includes:

{\def\LTcaptype{none} % do not increment counter
\begin{longtable}[]{@{}
  >{\raggedright\arraybackslash}p{(\linewidth - 2\tabcolsep) * \real{0.5000}}
  >{\raggedright\arraybackslash}p{(\linewidth - 2\tabcolsep) * \real{0.5000}}@{}}
\toprule\noalign{}
\begin{minipage}[b]{\linewidth}\raggedright
Property
\end{minipage} & \begin{minipage}[b]{\linewidth}\raggedright
Purpose
\end{minipage} \\
\midrule\noalign{}
\endhead
\bottomrule\noalign{}
\endlastfoot
\texttt{code} & Unique setting identifier and JSON/YAML path used for
object assignment. \\
\texttt{type} & Value type such as \texttt{boolean}, \texttt{integer},
\texttt{float}, \texttt{number}, \texttt{string}, \texttt{color}, or
\texttt{other}. \\
\texttt{default} & Reset value and fallback used when no saved value
exists. \\
\texttt{weight} & Sort order in the GUI and generated configuration
presentation. \\
\texttt{ingui} & Optional localization key for the GUI label. \\
\texttt{comment} & Optional literal help text or localization key. \\
\texttt{clientSide} & When \texttt{true}, keeps that setting client-side
instead of treating it as server-controlled. \\
\texttt{range} & Numeric \texttt{min}, \texttt{max}, and optional
\texttt{step}; a complete min/max range produces a slider. \\
\texttt{values} & Explicit allowed or suggested values for selection. \\
\texttt{mapping} & Maps display/config labels to underlying values. \\
\texttt{logarithmic} & Requests a logarithmic numeric slider. \\
\end{longtable}
}

Config Lib's decompiled enum also contains internal \texttt{None} and
\texttt{Constant} kinds. Constants are supplied separately through the
definition's \texttt{constants} object; they are useful to expressions
but are not ordinary editable controls.

\subsection{Formatting Blocks}\label{formatting-blocks}

Formatting blocks can appear among settings in the array form. For
example:

\begin{Shaded}
\begin{Highlighting}[]
\FunctionTok{\{}
  \DataTypeTok{"type"}\FunctionTok{:} \StringTok{"separator"}\FunctionTok{,}
  \DataTypeTok{"title"}\FunctionTok{:} \StringTok{"examplemod:settings{-}search{-}title"}\FunctionTok{,}
  \DataTypeTok{"weight"}\FunctionTok{:} \DecValTok{10}\FunctionTok{,}
  \DataTypeTok{"collapsible"}\FunctionTok{:} \KeywordTok{true}
\FunctionTok{\}}
\end{Highlighting}
\end{Shaded}

Separators can carry a \texttt{title} or explanatory \texttt{text}.
Config Lib attempts to resolve those strings as localization keys and
otherwise displays them literally.

\subsection{Grouped Form}\label{grouped-form}

The indexed upstream JSON API documents settings grouped by type:

\begin{Shaded}
\begin{Highlighting}[]
\FunctionTok{\{}
  \DataTypeTok{"version"}\FunctionTok{:} \DecValTok{1}\FunctionTok{,}
  \DataTypeTok{"settings"}\FunctionTok{:} \FunctionTok{\{}
    \DataTypeTok{"boolean"}\FunctionTok{:} \FunctionTok{\{}
      \DataTypeTok{"FEATURE\_ENABLED"}\FunctionTok{:} \FunctionTok{\{}
        \DataTypeTok{"name"}\FunctionTok{:} \StringTok{"feature{-}enabled"}\FunctionTok{,}
        \DataTypeTok{"ingui"}\FunctionTok{:} \StringTok{"examplemod:setting{-}feature{-}enabled"}\FunctionTok{,}
        \DataTypeTok{"comment"}\FunctionTok{:} \StringTok{"examplemod:setting{-}feature{-}enabled{-}description"}\FunctionTok{,}
        \DataTypeTok{"default"}\FunctionTok{:} \KeywordTok{true}\FunctionTok{,}
        \DataTypeTok{"weight"}\FunctionTok{:} \DecValTok{1}
      \FunctionTok{\}}
    \FunctionTok{\},}
    \DataTypeTok{"integer"}\FunctionTok{:} \FunctionTok{\{}
      \DataTypeTok{"SEARCH\_RADIUS"}\FunctionTok{:} \FunctionTok{\{}
        \DataTypeTok{"name"}\FunctionTok{:} \StringTok{"search{-}radius"}\FunctionTok{,}
        \DataTypeTok{"default"}\FunctionTok{:} \DecValTok{5}\FunctionTok{,}
        \DataTypeTok{"weight"}\FunctionTok{:} \DecValTok{2}\FunctionTok{,}
        \DataTypeTok{"range"}\FunctionTok{:} \FunctionTok{\{}
          \DataTypeTok{"min"}\FunctionTok{:} \DecValTok{1}\FunctionTok{,}
          \DataTypeTok{"max"}\FunctionTok{:} \DecValTok{20}\FunctionTok{,}
          \DataTypeTok{"step"}\FunctionTok{:} \DecValTok{1}
        \FunctionTok{\}}
      \FunctionTok{\}}
    \FunctionTok{\}}
  \FunctionTok{\}}
\FunctionTok{\}}
\end{Highlighting}
\end{Shaded}

Without \texttt{file}, this form generates
\texttt{ModConfig/examplemod.yaml}. The uppercase keys are setting codes
used by asset patches, while \texttt{name} is the user-facing YAML key.
The upstream documentation says YAML names should avoid spaces and
characters such as \texttt{:}; hyphens and underscores are suitable
separators.

\section{Asset Patching For Content
Mods}\label{asset-patching-for-content-mods}

Config Lib can patch asset values without any consuming-mod C\# code.
Patches are grouped by the value type they produce and then by asset
location:

\begin{Shaded}
\begin{Highlighting}[]
\FunctionTok{\{}
  \DataTypeTok{"version"}\FunctionTok{:} \DecValTok{1}\FunctionTok{,}
  \DataTypeTok{"settings"}\FunctionTok{:} \FunctionTok{\{}
    \DataTypeTok{"float"}\FunctionTok{:} \FunctionTok{\{}
      \DataTypeTok{"DAMAGE\_MULTIPLIER"}\FunctionTok{:} \FunctionTok{\{}
        \DataTypeTok{"name"}\FunctionTok{:} \StringTok{"damage{-}multiplier"}\FunctionTok{,}
        \DataTypeTok{"default"}\FunctionTok{:} \FloatTok{1.0}\FunctionTok{,}
        \DataTypeTok{"weight"}\FunctionTok{:} \DecValTok{1}\FunctionTok{,}
        \DataTypeTok{"range"}\FunctionTok{:} \FunctionTok{\{}
          \DataTypeTok{"min"}\FunctionTok{:} \FloatTok{0.5}\FunctionTok{,}
          \DataTypeTok{"max"}\FunctionTok{:} \FloatTok{2.0}\FunctionTok{,}
          \DataTypeTok{"step"}\FunctionTok{:} \FloatTok{0.1}
        \FunctionTok{\}}
      \FunctionTok{\}}
    \FunctionTok{\}}
  \FunctionTok{\},}
  \DataTypeTok{"patches"}\FunctionTok{:} \FunctionTok{\{}
    \DataTypeTok{"float"}\FunctionTok{:} \FunctionTok{\{}
      \DataTypeTok{"examplemod:itemtypes/example{-}tool.json"}\FunctionTok{:} \FunctionTok{\{}
        \DataTypeTok{"behaviors/0/properties/damage"}\FunctionTok{:} \StringTok{"value * DAMAGE\_MULTIPLIER"}
      \FunctionTok{\}}
    \FunctionTok{\}}
  \FunctionTok{\}}
\FunctionTok{\}}
\end{Highlighting}
\end{Shaded}

Here, \texttt{value} means the asset's original value at the patched
path. Config Lib supports setting references, arithmetic, constants such
as \texttt{pi} and \texttt{e}, and functions including trigonometric
operations, absolute value, square root, rounding, logarithms, min/max,
clamp, and conditional comparison helpers.

Patches can target the consuming mod's assets, vanilla assets, JSON
patch files, or another mod's assets. Cross-mod patching is useful for
server configuration add-ons, but the target paths and load-order
assumptions must be tested against the exact target versions.

\subsection{Conditional Patches}\label{conditional-patches}

The JSON API also supports boolean expressions that select between
values. A representative expression is:

\begin{Shaded}
\begin{Highlighting}[]
\StringTok{"some/path"}\ErrorTok{:} \StringTok{"(FEATURE\_ENABLED) ? NEW\_VALUE : value"}
\end{Highlighting}
\end{Shaded}

This keeps the original asset value when the feature is disabled.
Parenthesize nontrivial arithmetic and boolean expressions clearly;
compact expressions are easier to validate and diagnose from Config Lib
logs.

\section{Runtime API In Config Lib
1.12.0}\label{runtime-api-in-config-lib-1.12.0}

\subsection{\texorpdfstring{\texttt{IConfigProvider}}{IConfigProvider}}\label{iconfigprovider}

\texttt{ConfigLibModSystem} implements \texttt{IConfigProvider}. Its
relevant public surface includes:

\begin{itemize}
\tightlist
\item
  \texttt{Domains}, listing registered configuration domains;
\item
  \texttt{GetConfig(domain)} and \texttt{GetSetting(domain,\ code)};
\item
  \texttt{ConfigsLoaded}, \texttt{SettingChanged},
  \texttt{ConfigWindowOpened}, and \texttt{ConfigWindowClosed} events;
\item
  \texttt{RegisterCustomConfig}, for a custom-drawn configuration
  window;
\item
  \texttt{RegisterCustomManagedConfig}, for a configuration object
  managed by Config Lib.
\end{itemize}

\texttt{SettingChanged} supplies the domain, owning \texttt{IConfig},
and changed \texttt{ISetting}. Always filter the domain before modifying
a consuming mod's state.

\subsection{\texorpdfstring{\texttt{IConfig}}{IConfig}}\label{iconfig}

The public interface exposes:

\begin{itemize}
\tightlist
\item
  \texttt{ConfigFilePath} and \texttt{Version};
\item
  \texttt{WriteToFile()}, \texttt{ReadFromFile()}, and
  \texttt{RestoreToDefaults()};
\item
  \texttt{GetSetting(code)};
\item
  the per-config \texttt{SettingChanged} event;
\item
  \texttt{AssignSettingsValues(target)}.
\end{itemize}

Prefer this interface when no implementation-specific behavior is
needed. Casting to Config Lib's concrete \texttt{Config} class couples
the mod more tightly to the retained version.

\subsection{\texorpdfstring{\texttt{ISetting}}{ISetting}}\label{isetting}

The public setting interface exposes the current \texttt{Value},
optional \texttt{MappingKey}, \texttt{DefaultValue},
\texttt{SettingType}, \texttt{YamlCode}, validation metadata, and
\texttt{AssignSettingValue(target)}.

The assignment helper looks for a matching field or property on the
supplied object. Keep setting codes aligned with the configuration
object's actual member names, including capitalization and any
historical misspellings that have already become part of the saved-file
contract.

\subsection{Managed Registration}\label{managed-registration}

Config Lib \texttt{1.12.0} also exposes:

\begin{Shaded}
\begin{Highlighting}[]
\FunctionTok{RegisterCustomManagedConfig}\OperatorTok{(}
    \DataTypeTok{string}\NormalTok{ domain}\OperatorTok{,}
    \DataTypeTok{object}\NormalTok{ configObject}\OperatorTok{,}
    \DataTypeTok{string}\OperatorTok{?}\NormalTok{ path }\OperatorTok{=} \KeywordTok{null}\OperatorTok{,}
\NormalTok{    Action}\OperatorTok{?}\NormalTok{ onSyncedFromServer }\OperatorTok{=} \KeywordTok{null}\OperatorTok{,}
\NormalTok{    Action}\OperatorTok{\textless{}}\DataTypeTok{string}\OperatorTok{\textgreater{}?}\NormalTok{ onSettingChanged }\OperatorTok{=} \KeywordTok{null}\OperatorTok{,}
\NormalTok{    Action}\OperatorTok{?}\NormalTok{ onConfigSaved }\OperatorTok{=} \KeywordTok{null}\OperatorTok{)}
\end{Highlighting}
\end{Shaded}

It derives a definition from the supplied object's fields and
properties, uses \texttt{\textless{}domain\textgreater{}.json} when no
path is provided, assigns changed values back to the object, and offers
callbacks for synchronization, individual changes, and saves.

The implementation rejects registration after Config Lib has already
serialized its registry for clients. Register managed configurations
during early mod startup, not from late gameplay events. This is an
advanced alternative to the asset definition; the asset-driven JSON file
is easier to inspect, localize, and support without Config Lib code
references.

\section{Load Order And
Synchronization}\label{load-order-and-synchronization}

Config Lib \texttt{1.12.0} registers its configuration recipe registry
during \texttt{Start()} and loads discovered
\texttt{configlib-patches.json} assets during \texttt{AssetsLoaded()}.
After loading, it applies configured asset patches and raises
\texttt{ConfigsLoaded}.

The implementation distinguishes client-side settings from
server-controlled settings:

\begin{itemize}
\tightlist
\item
  \texttt{clientSide:\ true} marks a setting as local to the client;
\item
  server-controlled changes require the Vintage Story
  \texttt{controlserver} privilege;
\item
  server changes are sent through Config Lib's network channel and
  broadcast to other clients;
\item
  received settings are applied to the Config Lib configuration object
  before callbacks are raised.
\end{itemize}

Do not use \texttt{clientSide:\ true} for values that affect
authoritative gameplay, world state, recipes, damage, or server
simulation. Conversely, purely visual or local convenience settings
should not require server administration.

Config Lib also publishes event-bus names in \texttt{1.12.0}:

\begin{Shaded}
\begin{Highlighting}[]
\NormalTok{configlib:\textless{}domain\textgreater{}:config{-}saved}
\NormalTok{configlib:\textless{}domain\textgreater{}:setting{-}changed}
\NormalTok{configlib:\textless{}domain\textgreater{}:setting{-}loaded}
\NormalTok{configlib:config{-}reload}
\end{Highlighting}
\end{Shaded}

The strongly typed provider events are normally simpler for a code mod
that already references Config Lib. Event-bus integration can reduce
direct coupling but still requires exact payload handling for the target
Config Lib version.

\section{Optional Dependency Rules}\label{optional-dependency-rules}

The indexed upstream JSON documentation says Config Lib does not need to
be listed in \texttt{modinfo.json} for the asset-driven integration.
That is the cleanest option for content mods.

For code mods that call Config Lib types:

\begin{itemize}
\tightlist
\item
  compile against the intended Config Lib API;
\item
  do not call or resolve its mod system until
  \texttt{IsModEnabled("configlib")} succeeds;
\item
  isolate Config Lib-specific work in a guarded method;
\item
  keep all essential configuration loading outside that method;
\item
  verify a clean startup with Config Lib absent;
\item
  verify startup with Config Lib present but a malformed definition;
\item
  avoid packaging Config Lib's assembly as though it were part of the
  consuming mod.
\end{itemize}

Betterer Prospecting \texttt{3.4.5} demonstrates an optional
\texttt{modinfo.json} contract: its dependency list names only the game,
while its guarded runtime path checks for Config Lib before resolving
\texttt{ConfigLibModSystem}.

\section{Defaults, Versions, And
Migration}\label{defaults-versions-and-migration}

The most important maintenance rule is to keep defaults consistent
across all sources:

\begin{itemize}
\tightlist
\item
  C\# property or field initializers;
\item
  fallback asset values;
\item
  \texttt{configlib-patches.json} setting defaults;
\item
  any example configuration shipped in documentation.
\end{itemize}

Betterer Prospecting provides a concrete failure example. Its C\#
\texttt{ModConfig} initializes \texttt{BoreholeDmg} to \texttt{2}, while
its Config Lib manifest advertises \texttt{4}. The effective initial
behavior can therefore depend on which file was created and whether
Config Lib participated.

The upstream YAML documentation uses \texttt{version} to regenerate an
old generated configuration with new defaults and instructs authors to
increase it when defaults need to be regenerated. Treat that as a
migration boundary, not as a substitute for deliberate compatibility
handling. In particular:

\begin{itemize}
\tightlist
\item
  do not bump the version merely because the mod itself received a
  release;
\item
  preserve user choices where possible;
\item
  test the behavior of \texttt{version} separately for YAML and
  \texttt{file}-backed JSON definitions;
\item
  document intentional resets or renamed keys;
\item
  never silently correct a misspelled setting key after users may
  already have it in saved files.
\end{itemize}

\section{Localization}\label{localization}

For setting labels, Config Lib uses \texttt{ingui} when present and
otherwise falls back to the setting name/code. Comments can also be
translation keys. If a key has no domain prefix, Config Lib attempts to
use the consuming mod's asset domain; explicit
\texttt{\textless{}modid\textgreater{}:\textless{}key\textgreater{}}
values are clearer in shared examples.

Add matching strings to the mod's language assets, for example:

\begin{Shaded}
\begin{Highlighting}[]
\FunctionTok{\{}
  \DataTypeTok{"setting{-}enabled"}\FunctionTok{:} \StringTok{"Enable feature"}\FunctionTok{,}
  \DataTypeTok{"setting{-}enabled{-}description"}\FunctionTok{:} \StringTok{"Turns the example feature on or off."}\FunctionTok{,}
  \DataTypeTok{"setting{-}search{-}radius"}\FunctionTok{:} \StringTok{"Search radius"}\FunctionTok{,}
  \DataTypeTok{"setting{-}search{-}radius{-}description"}\FunctionTok{:} \StringTok{"Maximum horizontal search distance."}
\FunctionTok{\}}
\end{Highlighting}
\end{Shaded}

Test both translated and missing-key behavior. Config Lib falls back to
displaying literal text when a separator or description cannot be
resolved as a translation.

\section{Recommended Verification
Matrix}\label{recommended-verification-matrix}

Before release, test at least the following:

{\def\LTcaptype{none} % do not increment counter
\begin{longtable}[]{@{}
  >{\raggedright\arraybackslash}p{(\linewidth - 2\tabcolsep) * \real{0.5000}}
  >{\raggedright\arraybackslash}p{(\linewidth - 2\tabcolsep) * \real{0.5000}}@{}}
\toprule\noalign{}
\begin{minipage}[b]{\linewidth}\raggedright
Scenario
\end{minipage} & \begin{minipage}[b]{\linewidth}\raggedright
Expected Result
\end{minipage} \\
\midrule\noalign{}
\endhead
\bottomrule\noalign{}
\endlastfoot
Config Lib absent & Mod loads; normal JSON/default configuration works;
no Config Lib errors occur. \\
Config Lib present, first run & Config appears in the GUI and the
expected JSON/YAML file is generated or updated. \\
Existing valid config & Saved values survive restart and appear
correctly in the GUI. \\
Invalid or partial config & Mod falls back safely or reports a useful
error without losing unrelated functionality. \\
Live setting edit & Matching config member changes and dependent runtime
state refreshes. \\
Dedicated server edit & Only authorized server settings change; clients
receive the synchronized values. \\
Client-only edit & Local visual/convenience behavior changes without
modifying server-authoritative state. \\
Reset to defaults & C\# defaults, asset defaults, and Config Lib
defaults agree. \\
Config Lib removed after use & The mod still reads its ordinary config
or uses sensible fallback defaults. \\
\end{longtable}
}

For asset patches, also inspect the target asset after loading and test
with every targeted mod/version combination. An asset path that worked
in one game patch or dependency release is not a permanent API.

\section{Known Limitations And
Uncertainty}\label{known-limitations-and-uncertainty}

\begin{itemize}
\tightlist
\item
  The upstream JSON-API page indexed here was last known edited on
  \texttt{2024-05-14} and retrieved on \texttt{2026-08-10}. Config Lib
  \texttt{1.12.0} contains newer behavior, including array-form
  settings, JSON \texttt{file} backing, color settings, and managed
  object registration that is not fully explained on that page.
\item
  The retained Config Lib release is a same-line fallback for Vintage
  Story \texttt{1.22.6}, not an exact patch-tag match.
\item
  Decompiled source is behavioral evidence, not build-ready source and
  not a substitute for the upstream public package/API documentation.
\item
  Config Lib's asset patch system is intentionally narrower than Vintage
  Story's complete JSON patch feature set.
\item
  Optional direct C\# linkage should be tested in the consuming mod's
  actual build and packaging arrangement. The Betterer Prospecting
  pattern demonstrates that it can work, but it does not prove every
  assembly layout is safe.
\item
  The precise GUI capabilities depend on Config Lib and its VSImGui
  dependency being compatible with the client.
\end{itemize}

\section{Evidence And Source Paths}\label{evidence-and-source-paths}

The research used the Vintage Story Reference \texttt{1.22.6} mod
manifest, web index, asset index, and code index before opening the
following narrow source candidates.

\subsection{Indexed Upstream
Documentation}\label{indexed-upstream-documentation}

\begin{itemize}
\tightlist
\item
  Config Lib JSON API:
  \url{https://github.com/maltiez2/vsmod_configlib/wiki/JSON-API}
\item
  Indexed raw Markdown endpoint:
  \url{https://raw.githubusercontent.com/wiki/maltiez2/vsmod_configlib/JSON-API.md}
\item
  Indexed retrieval time: \texttt{2026-08-10T03:39:34.076833+00:00}
\item
  Last-known upstream edit recorded by the source manifest:
  \texttt{2024-05-14}
\end{itemize}

\subsection{Config Lib 1.12.0
Reference}\label{config-lib-1.12.0-reference}

All paths begin at \texttt{versions/1.22.6/mods/configlib\_1.12.0/}.

\begin{itemize}
\tightlist
\item
  \texttt{modinfo.json} --- retained package identity, VSImGui
  dependency, and side metadata.
\item
  \texttt{decompiled\ assembly/configlib/ConfigLib/ConfigLibModSystem.cs}
  --- asset discovery, lifecycle, networking, public provider API,
  managed registration, and event-bus constants.
\item
  \texttt{decompiled\ assembly/configlib/ConfigLib/Config.cs} ---
  YAML/JSON backing, version parsing, setting assignment, formatting,
  constants, and file read/write behavior.
\item
  \texttt{decompiled\ assembly/configlib/ConfigLib/ConfigPatches.cs} ---
  collection and application of asset patches.
\item
  \texttt{decompiled\ assembly/configlib/ConfigLib/ConfigSetting.cs} ---
  setting parsing, validation, mapping, side metadata, and
  reflection-based assignment.
\item
  \texttt{decompiled\ assembly/configlib/ConfigLib/IConfigProvider.cs}
  --- public provider contract.
\item
  \texttt{decompiled\ assembly/configlib/ConfigLib/IConfig.cs} ---
  public configuration contract.
\item
  \texttt{decompiled\ assembly/configlib/ConfigLib/ISetting.cs} ---
  public setting contract.
\item
  \texttt{decompiled\ assembly/configlib/ConfigLib/ConfigSettingType.cs}
  --- setting type enum.
\end{itemize}

\subsection{Real Consumer: Betterer Prospecting
3.4.5}\label{real-consumer-betterer-prospecting-3.4.5}

All paths begin at
\texttt{versions/1.22.6/mods/bettererprospecting\_3.4.5/}.

\begin{itemize}
\tightlist
\item
  \texttt{assets/bettererprospecting/config/configlib-patches.json} ---
  array-form settings, JSON \texttt{file} backing, separators,
  localization keys, ranges, client-side setting, and version
  \texttt{31}.
\item
  \texttt{decompiled\ assembly/BetterErProspecting/BetterErProspecting/BetterErProspect.cs}
  --- normal config loading, optional Config Lib guard, provider events,
  object assignment, and runtime refresh behavior.
\item
  \texttt{decompiled\ assembly/BetterErProspecting/BetterErProspecting/Config/ModConfig.cs}
  --- ordinary fallback config and the \texttt{BoreholeDmg} default
  mismatch.
\item
  \texttt{modinfo.json} --- no hard Config Lib dependency declaration.
\end{itemize}

\subsection{Representative Index
Queries}\label{representative-index-queries}

\begin{Shaded}
\begin{Highlighting}[]
\NormalTok{python }\OperatorTok{{-}}\NormalTok{B tools}\OperatorTok{/}\NormalTok{reference{-}query}\OperatorTok{.}\FunctionTok{py} \OperatorTok{{-}{-}}\NormalTok{version }\DecValTok{1.22}\OperatorTok{.}\DecValTok{6}\NormalTok{ \textasciigrave{}}
\NormalTok{  mod }\OperatorTok{{-}{-}}\NormalTok{mod{-}id configlib }\OperatorTok{{-}{-}}\NormalTok{format json}

\NormalTok{python }\OperatorTok{{-}}\NormalTok{B tools}\OperatorTok{/}\NormalTok{reference{-}query}\OperatorTok{.}\FunctionTok{py} \OperatorTok{{-}{-}}\NormalTok{version }\DecValTok{1.22}\OperatorTok{.}\DecValTok{6}\NormalTok{ \textasciigrave{}}
\NormalTok{  web }\OperatorTok{{-}{-}}\NormalTok{source configlib{-}json{-}api }\OperatorTok{{-}{-}}\NormalTok{limit }\DecValTok{100}

\DataTypeTok{python} \OperatorTok{{-}}\NormalTok{B tools}\OperatorTok{/}\NormalTok{reference{-}query}\OperatorTok{.}\FunctionTok{py} \OperatorTok{{-}{-}}\NormalTok{version }\DecValTok{1.22}\OperatorTok{.}\DecValTok{6}\NormalTok{ \textasciigrave{}}
\NormalTok{  asset }\OperatorTok{{-}{-}}\NormalTok{search configlib }\OperatorTok{{-}{-}}\NormalTok{limit }\DecValTok{30}

\DataTypeTok{python} \OperatorTok{{-}}\NormalTok{B tools}\OperatorTok{/}\NormalTok{reference{-}query}\OperatorTok{.}\FunctionTok{py} \OperatorTok{{-}{-}}\NormalTok{version }\DecValTok{1.22}\OperatorTok{.}\DecValTok{6}\NormalTok{ \textasciigrave{}}
\NormalTok{  code }\OperatorTok{{-}{-}}\NormalTok{root }\StringTok{"configlib\_1.12.0/decompiled assembly"}\NormalTok{ \textasciigrave{}}
  \OperatorTok{{-}{-}}\NormalTok{type IConfigProvider }\OperatorTok{{-}{-}}\NormalTok{limit }\DecValTok{10}

\DataTypeTok{python} \OperatorTok{{-}}\NormalTok{B tools}\OperatorTok{/}\NormalTok{reference{-}query}\OperatorTok{.}\FunctionTok{py} \OperatorTok{{-}{-}}\NormalTok{version }\DecValTok{1.22}\OperatorTok{.}\DecValTok{6}\NormalTok{ \textasciigrave{}}
\NormalTok{  code }\OperatorTok{{-}{-}}\NormalTok{root }\StringTok{"configlib\_1.12.0/decompiled assembly"}\NormalTok{ \textasciigrave{}}
  \OperatorTok{{-}{-}}\NormalTok{member AssignSettingValue }\OperatorTok{{-}{-}}\NormalTok{limit }\DecValTok{10}
\end{Highlighting}
\end{Shaded}

\section{Implementation Checklist}\label{implementation-checklist}

\begin{itemize}
\tightlist
\item[$\square$]
  Normal config loading works without Config Lib.
\item[$\square$]
  \texttt{configlib-patches.json} is under the exact domain/config path.
\item[$\square$]
  \texttt{file} matches the code mod's exact JSON filename when JSON
  backing is intended.
\item[$\square$]
  Setting codes match config member names or intended JSON paths.
\item[$\square$]
  Defaults agree across C\#, assets, Config Lib, and documentation.
\item[$\square$]
  Server-authoritative and \texttt{clientSide} settings are classified
  correctly.
\item[$\square$]
  Config Lib calls are guarded by \texttt{IsModEnabled("configlib")}.
\item[$\square$]
  Initial values are assigned after Config Lib reports configs loaded.
\item[$\square$]
  Live callbacks refresh any derived runtime state.
\item[$\square$]
  Event handlers are removed during disposal when the mod system can be
  reloaded.
\item[$\square$]
  Localization keys exist and missing-key fallbacks were inspected.
\item[$\square$]
  Config Lib absent/present, client/server, restart, reset, and
  migration cases were tested.
\end{itemize}

\end{document}
